\documentclass[12pt]{book}
\usepackage[utf8]{inputenc}
\usepackage[T1]{fontenc}
\usepackage{amsmath,amssymb,amsfonts}
\usepackage{mathrsfs}
\usepackage{amsthm}

% Calligraphic for categories and sheaves
\newcommand{\cat}[1]{\mathcal{#1}}
\newcommand{\sheaf}[1]{\mathcal{#1}}

% Standard math operators
\DeclareMathOperator{\Hom}{Hom}
\DeclareMathOperator{\Ext}{Ext}
\DeclareMathOperator{\Tor}{Tor}
\DeclareMathOperator{\id}{id}
\DeclareMathOperator{\Rfst}{\mathbf{R}f_*}
\DeclareMathOperator{\Lfst}{\mathbf{L}f_*}
\DeclareMathOperator{\RHom}{\mathbf{R}\mathcal{H}\textit{om}}

\begin{document}

\setcounter{page}{266}
\section{The Uniqueness of the Dualizing Complex.}

\begin{theorem}
Let \(X\) be a connected, locally noetherian prescheme, and let \(R^{\bullet}\) be a dualizing complex on \(X\).
Let \(R'^{\bullet} \in D_{C}^{+}(X)\) be any other complex.
Then \(R'^{\bullet}\) is a dualizing complex if and only if there exists an invertible sheaf \(L\) on \(X\), and an integer \(n\), such that
\[
R'^{\bullet} \cong R^{\bullet} \otimes L[n] .
\]
Furthermore, \(L\) and \(n\) are determined uniquely to within isomorphism.
\end{theorem}

\begin{proof}
1) Observe that such an \(R'^{\bullet}\) is a dualizing complex. Indeed, we may assume that \(R^{\bullet}\) is a bounded complex of quasicoherent injectives. Then \(R^{\bullet} \otimes L[n]\) is also a bounded complex of injectives, since the property of being injective is local [II 7.16], and \(L\) is locally free. \(R'\) is dualizing by Corollary 2.3.

2) Observe that \(L\) and \(n\) are uniquely determined by the identity
\[
L[n] \cong \RHom^{\cdot}(R^{\cdot}, R'^{\cdot}) .
\]

\setcounter{page}{267}
3) Now we show that existence of \(L\) and \(n\), assuming that \(R'\) is a dualizing complex.
Let \(D\) and \(D'\) be the dualizing functors corresponding to \(R^{\bullet}\) and \(R'^{\bullet}\).
Define
\[
L^{\cdot}=D' D(\sheaf{X})=\RHom^{\cdot}(R^{\cdot}, R'^{\cdot}) .
\]
We will prove that \(L^{\cdot}\) is isomorphic, in \(D^{+}(X)\), to a complex of the form \(L[n]\), where \(L\) is an invertible sheaf, and that
\[
R'^{\cdot} \cong R^{\cdot} \otimes L^{\cdot}.
\]
\end{proof}

\begin{lemma}
There is a natural functorial isomorphism, for all \(M^{\cdot} \in D_{C}^{-}(X)\),
\[
M^{\cdot} \otimes L^{\cdot} \stackrel{\sim}{\to} D' D\left(M^{\cdot}\right) .
\]
(Note that this tensor product makes sense since \(L^{\cdot} \in D_{C}^{-}(X)\), cf. [II ~].)
\end{lemma}

\begin{proof}
AS in Lemma 1.2 above, one defines a functorial homomorphism
\[
M^{*} \otimes \RHom^{\cdot}\left(R^{*}, R'^{\cdot}\right)
\to \RHom^{\cdot}\left(\RHom^{\cdot}\left(M^{*}, R^{*}\right), R'^{\cdot}\right) .
\]
Since it is an isomorphism for \(M^{\cdot}=\sheaf{X}\) by the lemma on way-out functors, it is an isomorphism for all \(M^{*} \in D_{C}^{-}(X)\).

\setcounter{page}{268}
We now define \(L'^{\cdot}=DD'(\sheaf{X})\), and use the lemma to deduce that
\[
L^{\cdot} \otimes L'^{\cdot}=D' D\left(L'^{\cdot}\right)=D' D D'\left(\sheaf{X}\right)=\sheaf{X} .
\]
\end{proof}

\begin{lemma}
Let \(X\) be a connected locally noetherian prescheme, and let \(L^{\cdot}, L'^{\cdot} \in D_{C}^{-}(X)\) be such that
\[
L^{\cdot} \otimes L'^{\cdot} \cong \sheaf{X}.
\]
Then \(L^{\cdot}=L[n]\) for some invertible sheaf \(L\), and \(L'^{\cdot}=L^{\vee}[-n]\).
\end{lemma}

\begin{proof}
Since \(X\) is connected, we reduce immediately to the case where \(X\) is the spectrum of a local noetherian ring \(A\), and we wish to show that \(L^{\cdot}=\sheaf{X}[n]\).
Let \(p\) be the largest integer such that \(H^{p}(L^{\cdot}) \ne 0\), and let \(q\) be the largest integer such that \(H^{q}(L'^{\cdot}) \ne 0\).
Then one sees easily that \(p+q\) is the largest integer for which \(H^{n}(L^{\cdot} \otimes L'^{\cdot}) \ne 0\), and
\[
H^{p+q}(L^{\cdot} \otimes L'^{\cdot})
=H^{p}(L^{\cdot}) \otimes H^{q}(L'^{\cdot}) .
\]
Therefore \(p+q=0\), and
\[
H^{p}(L^{\cdot}) \otimes H^{q}(L'^{\cdot})=\sheaf{X}.
\]
It follows, by a simple lemma on modules of finite type over a local noetherian ring [cf. EGA, Ch. O, 5.4.3],
that \(H^{p}(L^{\cdot}) \cong H^{q}(L'^{\cdot}) \cong \sheaf{X}\).
Now since \(\sheaf{X}\) is locally free, one shows that \(L^{\cdot}=L_{i} \oplus \sheaf{X}[-p]\),
where \(L_{i}\) has cohomology only in dimensions \(< p\).
Similarly, \(L'^{\cdot}=L_{1}' \oplus \sheaf{X}[p]\).

\setcounter{page}{269}
Thus from
\[
L^{\cdot} \otimes L'^{\cdot} \cong \sheaf{X}
\]
we deduce
\[
L_{i}[p] \oplus L_{1}'[-p] \oplus\left(L_{1}^{\cdot} \otimes L_{1}'^{\cdot}\right)=0
\]
from which it follows that \(L_{i}\) and \(L_{1}'\) are \(0\) in \(D_{c}^{-}(X)\), so that \(L^{\cdot} \cong \sheaf{X}[-p]\),
and \(L'^{\cdot} \cong \sheaf{X}[p]\), as required.
\end{lemma}

This proves the lemma, so we have only to show that
\[
R'^{\cdot}=R^{\cdot} \otimes L^{\cdot}.
\]
Indeed, \(R^{\cdot} \in D_{c}^{b}(X)\), so applying Lemma 3.2,
\[
R^{*} \otimes L^{*}=D' D\left(R^{*}\right)=D'\left(\sheaf{X}\right)=R'^{\cdot} \quad q.e.d.
\]

We now give an application of the uniqueness theorem.

\begin{proposition}
Let \(A\) be a noetherian local ring with residue field \(k\), and let \(R^{\cdot} \in D_{C}^{+}(A)\).
Then \(R^{\cdot}\) is dualizing if and only if there is an integer \(d\) such that
\[
Ext^{i}\left(k, R^{\cdot}\right)=
\begin{cases}
0 & i \neq d \\
k & i=d .
\end{cases}
\]
(Here we write \(D^{+}(A)\) for the derived category of the category of \(A\)-modules.
The subscript \(c\) denotes complexes whose cohomology modules are of finite type.
We carry over the definitions and results of the previous sections to this case.
(Cf. [II.7.19] which ensures that we will not get into trouble.))
\end{proposition}

\begin{proof}
First suppose that \(R^{\cdot}\) is dualizing, and let \(j : \operatorname{Spec}k \to \operatorname{Spec}A\) be the inclusion.
Then by Proposition 2.4,

\setcounter{page}{270}
\[
j^{b}\left(R^{\cdot}\right)=\RHom^{\cdot}\left(k, R^{\cdot}\right)
\]
is a dualizing complex on \(k\).
But \(k\) itself is a dualizing complex on \(k\), so by the Uniqueness Theorem above,
\(j^{b}(R^{\cdot}) \cong k[-d]\) for some \(d\). This \(d\) will do.

For the converse, it is clear that \(k\) is reflexive, so by Proposition 2.5 we have only to show that \(R^{\cdot}\) has finite inJective dimension.
Indeed, we will show by induction on \(r=\dim M\), for any \(A\)-module \(M\) of finite type, that
\[
Ext^{i}\left(M, R^{\cdot}\right)=0 \text{ for } i \notin[d-r, d].
\]
Indeed, for \(\dim M=0\), \(M\) is of finite length, and our statement follows from the case \(M=k\) by induction on the length of \(M\).

For \(\dim M=r>0\), we induct as in the proof of Proposition 2.5.
First we may assume that \(\mathfrak{m} \notin \operatorname{Ass}M\), where \(\mathfrak{m}\) is the maximal ideal of \(A\).
Then choosing \(t \in \mathfrak{m}\) with \(t\) a non-zerodivisor in \(M\), we have
\[
0 \to M \stackrel{t}{\to} M \to M / tM \to 0 .
\]
Then for \(i \notin[d-r+1, d]\) we have \(Ext^{i}(M / t M, R^{*})=0\),
hence for \(i \notin[d-r,d]\) we have from the long exact sequence of Ext's,
\[
Ext^{i}\left(M, R^{\cdot}\right) \stackrel{t}{\to} Ext^{i}\left(M, R^{\cdot}\right) \to 0 .
\]
It follows by Nakayama's lemma that \(Ext^{i}(M, R^{\cdot})=0\). q.e.d.
\end{proof}

\setcounter{page}{271}
\begin{corollary}
Let \(A\) be a noetherian local ring, and \(R^{*} \in D_{C}^{+}(A)\).
Then \(R^{*}\) is dualizing on \(A\) if and only \(R \otimes_{A} \hat{A}\) is dualizing on the completion \(\hat{A}\) of \(A\).
\end{corollary}

\begin{proof}
clearly
\[
R^{\cdot} \otimes_{A} \hat{A} \in D_{C}^{+}(\hat{A}) .
\]
Furthermore,
\[
Ext_{\hat{A}}^{i}\left(k, R^{\cdot} \otimes_{A} \hat{A}\right)
=Ext_{A}^{i}\left(k, R^{\cdot}\right) \otimes_{A} \hat{A}
\]
for all \(i\). Thus the "only if" implication is clear.
For the "if" implication, note that the functor \(\otimes_{A} \hat{A}\) is faithfully flat,
so if \(M\) is a non-zero \(A\)-module, then \(M \otimes_{A} \hat{A}\) is non-zero.
Furthermore, if \(M\) is an \(A\)-module, such that \(M \otimes_{A} \hat{A} \cong k\), then \(M \cong k\).
(Indeed, map \(M\) to \(k\) via the natural map \(M \to M \otimes_{A} \hat{A}\).
This map becomes an isomorphism upon tensoring with \(\hat{A}\), hence was an isomorphism).
The result thus follows from the proposition.
\end{corollary}

\end{document}